% ============================================================================
% บทที่ 7 — เก็บงาน เปิดงานต่อ และดาวน์โหลดผลลัพธ์
% ============================================================================
\mychapterlong{เก็บงาน เปิดงานต่อ และดาวน์โหลดผลลัพธ์}{เก็บงาน เปิดงานต่อ\\ และดาวน์โหลดผลลัพธ์}

\noindent บทนี้ตอบคำถามเดียว: จะเก็บงานนี้ เปิดทำต่อ หรือส่งต่อให้ผู้อื่นอย่างไร

\section{วิธีเก็บงานและเปิดงานต่อ}

\begin{center}
\footnotesize
\begin{tabular}{L{2.5cm}ccccL{2.7cm}}
\toprule
\textbf{วิธี} & \textbf{เปิดต่อในระบบ} & \textbf{ข้ามเครื่อง} & \textbf{เก็บการแก้ไข} & \textbf{เปิดดูภายนอก} & \textbf{เหมาะกับ} \\
\midrule
สถานะในเบราว์เซอร์ & \yesmark & \nomark & \yesmark & \nomark & ทำงานต่อบนเครื่องเดิม \\
บันทึกบนคลาวด์ & \yesmark & \yesmark & \yesmark & ลิงก์ & ทำงานต่อ/แชร์ \\
\file{project.zip} & \nomark & ไฟล์ส่งต่อได้ & สำเนาผลลัพธ์ & \yesmark & เก็บถาวร/ตรวจสอบ \\
นำเข้าไฟล์ต้นทางใหม่ & \yesmark & \yesmark & \nomark & \nomark & เริ่มจากต้นทางใหม่ \\
\bottomrule
\end{tabular}
\end{center}

\begin{caution}
\textbf{ไฟล์ \file{project.zip} นำกลับมาเปิดงานต่อในระบบไม่ได้}
ระบบไม่มีฟังก์ชันนำเข้าไฟล์นี้ ไฟล์ \file{project.zip} ใช้เปิดดู ตรวจสอบ
หรือส่งต่อผลลัพธ์เท่านั้น การเปิดงานต่อต้องใช้การกู้คืนในเบราว์เซอร์เดิม
การเปิดลิงก์แชร์ หรือการนำเข้าไฟล์ต้นทางใหม่
\end{caution}

การนำเข้าไฟล์ต้นทางใหม่จะเริ่มโครงการใหม่ทั้งหมด จึงไม่นำการล็อกเวลา
ที่ผู้ใช้เคยแก้ไว้กลับมา ส่วนการจัดตารางใหม่ (ไม่ใช่นำเข้าใหม่) มีผลต่อ
การล็อกเวลาดังที่อธิบายไว้ในบทที่ 5

\section{การดาวน์โหลดไฟล์เก็บถาวร}

\begin{steps}
  \item ไปที่เมนู \textbf{ดาวน์โหลด project.zip} ที่ด้านซ้ายล่าง
\end{steps}

ได้ไฟล์ \file{project.zip} ซึ่งเป็นไฟล์เก็บถาวรผลลัพธ์ของโครงการ
ประกอบด้วยข้อมูลโครงการ รายงานตรวจสอบ และไฟล์ส่งออกในโฟลเดอร์
\texttt{exports/} (CSV และ HTML) ไฟล์ทั้งหมดในไฟล์เก็บถาวรอยู่ในภาคผนวก B
ไฟล์ที่ผู้ใช้งานประจำเปิดบ่อยที่สุดคือ:

\begin{center}
\small
\begin{tabular}{L{6.5cm}L{8cm}}
\toprule
\textbf{ไฟล์ (ในโฟลเดอร์ exports/)} & \textbf{เนื้อหา} \\
\midrule
ตารางสอบทั้งหมด.csv & ตารางสอบรวมทุกภาค \\
ตารางสอบกลางภาค.csv & ตารางสอบกลางภาค \\
ตารางสอบปลายภาค.csv & ตารางสอบปลายภาค \\
ปฏิทินสอบกลางภาค.html & ปฏิทินกลางภาคแบบเปิดในเบราว์เซอร์ \\
ปฏิทินสอบปลายภาค.html & ปฏิทินปลายภาคแบบเปิดในเบราว์เซอร์ \\
\bottomrule
\end{tabular}
\end{center}

\begin{caution}
ไฟล์เก็บถาวรจากเบราว์เซอร์ไม่รวมไฟล์ข้อมูลต้นทาง เช่น
\file{managed.xlsx} หรือ \file{context.xlsx} ดังนั้นควรเก็บไฟล์ต้นฉบับ
แยกไว้เองเสมอ
\end{caution}

\section{การเปิดไฟล์ตารางสอบ CSV}

\begin{steps}
  \item แตกไฟล์ \file{project.zip} ที่ดาวน์โหลดมา
  \item เปิดโฟลเดอร์ \texttt{exports}
  \item เปิดไฟล์ \file{ตารางสอบทั้งหมด.csv} ด้วยโปรแกรมตารางคำนวณ
\end{steps}

ตาราง CSV มีคอลัมน์ วันที่ เวลา รหัสวิชา ชื่อวิชา ตอนเรียน กลุ่มนักศึกษา
ที่มา และสถานะ โดยคอลัมน์เวลาใช้รูปแบบกระชับ
เช่น \texttt{9-12} และ \texttt{13:30-16:30} ตัวอย่างเนื้อหาไฟล์:

\begin{center}
\small
\begin{tabular}{L{2.1cm}L{0.9cm}L{2.1cm}L{2.5cm}L{0.9cm}L{1.6cm}L{1.7cm}L{1.3cm}}
\toprule
\textbf{วันที่} & \textbf{เวลา} & \textbf{รหัสวิชา} & \textbf{ชื่อวิชา} & \textbf{ตอนเรียน} & \textbf{กลุ่มนักศึกษา} & \textbf{ที่มา} & \textbf{สถานะ} \\
\midrule
2026-08-17 & 9-12 & \course{061000003} & ระบบฐานข้อมูลเบื้องต้น & 1 & DEMO-C & สร้างอัตโนมัติ & จัดแล้ว \\
2026-08-17 & 9-12 & \course{061000004} & การวิเคราะห์และออกแบบระบบ & 1 & DEMO-D & สร้างอัตโนมัติ & จัดแล้ว \\
\bottomrule
\end{tabular}
\end{center}

ค่าคอลัมน์ \textbf{ที่มา} เป็น \texttt{สร้างอัตโนมัติ} เมื่อระบบจัดเวลาเอง
\texttt{กำหนดเอง} เมื่อผู้ใช้แก้ไขเอง และ \texttt{นำเข้า} เมื่อเวลามาจากไฟล์ต้นทาง

\manualfigure{docs/usage-report/screenshots/crops/F24a.png}{0.95\textwidth}{%
  ไฟล์ \file{ตารางสอบทั้งหมด.csv} ที่เปิดในโปรแกรมตารางคำนวณ
  (แสดงเฉพาะส่วนต้นของไฟล์)}

\begin{note}
การแก้ไขไฟล์ CSV ที่ดาวน์โหลดมาไม่ได้เปลี่ยนตารางสอบในระบบ
ไฟล์ CSV เป็นเพียงผลลัพธ์ที่ส่งออกเท่านั้น
\end{note}

\section{การเปิดปฏิทิน HTML}

\begin{steps}
  \item เปิดไฟล์ \file{ปฏิทินสอบกลางภาค.html} หรือ \file{ปฏิทินสอบปลายภาค.html}
        ด้วยเบราว์เซอร์
\end{steps}

\manualfigure{docs/usage-report/screenshots/crops/F24b.png}{0.8\textwidth}{%
  ปฏิทินสอบกลางภาคในรูปแบบ HTML}

\section{การดาวน์โหลดรายงานตรวจสอบ}

หากผู้ดูแลระบบหรือฝ่ายสนับสนุนขอข้อมูลสำหรับตรวจสอบ
สามารถดาวน์โหลดรายงานตรวจสอบจากหน้า \textbf{ตรวจสอบ} ได้
รายละเอียดของไฟล์รายงานอยู่ในภาคผนวก B

\section{การกู้คืนสถานะในเบราว์เซอร์}

เมื่อกลับมาเปิดระบบด้วยเบราว์เซอร์และโปรไฟล์เดิม
ระบบจะกู้คืนสถานะล่าสุดที่เก็บไว้ในเครื่องโดยอัตโนมัติ
ทั้งข้อมูลโครงการที่นำเข้าและผลลัพธ์ที่แก้ไขไว้

\begin{caution}
การกู้คืนสถานะใช้ได้กับเบราว์เซอร์และโปรไฟล์เดิมเท่านั้น
ไม่ครอบคลุมการกู้คืนจากอุปกรณ์อื่น
\end{caution}

\manualfigure{docs/usage-report/screenshots/figures/F25.png}{0.95\textwidth}{%
  สถานะที่กู้คืนหลังเปิดเบราว์เซอร์ใหม่}

\seealso{บทที่ 8 สำหรับการบันทึกบนคลาวด์และการแชร์ และบทที่ 5 สำหรับการแก้ไขเวลา}
